%% kapa: exam=1 round=29 year=2018 period=1 form=A subject=civil_law subject_ko=민법 seq=27 number=27
%% source: Q-Net 감정평가사 기출문제 (cst00309) artlSeq=5202808 file=extracted/1008/5202808/2018년 제29회 감정평가사 1차시험 기출문제/2018년 제29회 감정평가사 1차시험 기출문제/1교시 A형.pdf
%% conversion: kordoc 4.24.51 -> tools/gen1.py (auto); review=proofread-sonnet; flags=none
\begin{question}{27}
다음 설명 중 옳은 것은? (다툼이 있으면 판례에 따름)
\begin{choices}
  \item 미등기 무허가건물의 양수인은 소유권에 준하는 관습상의 물권을 취득한다.
  \item 등기는 물권의 존속요건이므로, 등기가 불법 말소되면 물권은 소멸한다.
  \item 소유권이전등기의 원인으로 주장된 계약서가 진정하지 않은 것으로 증명되어도 그 등기의 적법추정은 복멸되지 않는다.
  \item 지하 또는 지상의 공간은 상하의 범위를 정하여 건물 기타 공작물을 소유하기 위한 구분지상권의 목적으로 할 수 없다.
  \item 공유자 중 1인이 다른 공유자의 동의 없이 그 공유 토지의 특정부분을 매도하여 타인명의로 소유권이전등기를 마친 경우, 그 매도부분 토지에 관한 소유권이전등기는 처분공유자의 공유지분 범위 내에서는 유효한 등기이다.
\end{choices}
\end{question}
\answer{5}
